% results_rq5

\section{RQ5: Feature Bloat and Companion Hardware}
\label{sec:rq5}

A fourth framing holds that adding features drives complaints about clutter and
complexity --- that a leaner app is a more satisfying one. We tested this on a
pre-specified hypothesis and test, before running the full corpus
(\S\ref{sec:rq5-provenance}). It fails in both directions: feature count does
not predict bloat complaints, and if anything the relationship runs backward.
What users mean by ``bloat'' is a property of the interface, not a count of
capabilities.

\subsection{Feature Count Does Not Predict Bloat Complaints}

The primary mixed-effects specification (M4) did not converge on its
\texttt{feature\_count} term (\(\beta = -0.004\), SE \(743.8\),
\(p \approx 0.99999\)), so we report this null from the cluster-robust
sensitivity fit, which did: the interaction
\texttt{complaint\_complexity\_bloat} \(\times\) \texttt{feature\_count} gives
\(\beta = -0.000487\) (95\% CI \([-0.00136, 0.00038]\), \(p = 0.272\)),
clustered on 26 apps. Adding features does not make users more likely to call
an app bloated. Nor is the null an artefact of a few large, well-known apps
absorbing the variance: feature count correlates only weakly with install count
across the 26 (Spearman \(\rho = 0.244\)).

\subsection{If Anything, the Relationship Runs Backward}

Figure~\ref{fig:rq5-bloat-curve} shows bloat-complaint rate against feature
count, and the pattern is the opposite of what the bloat hypothesis predicts.
The five leanest apps (6--9 features) average 3.96 stars at a bloat-complaint
rate of 5.2\%; the five fullest (13--14 features) average 4.71 stars at 1.6\%
--- a 3.3\(\times\) lower rate in the apps with the \emph{most} features,
against a corpus mean of 4.40 stars. The lean group is not uniquely defined:
six applications tie at nine features, so the bottom five can be drawn fifteen
ways, across which its mean rating spans 3.88 to 4.14. We report the set our
selection rule returned and treat the ratio as indicative rather than exact;
the fullest group carries no such tie. The pattern survives dropping the two
smallest apps by review volume (\(n=5\) and \(n=45\)): the remaining lean apps
still average 4.20 stars at 5.6\%. The clearest pair is iPray, 8 features and
2,225 reviews, at 8.0\%, against Athan, 13 features and 44,954 reviews, at
1.5\%. An exploratory app-level correlation, entered into the correction and
surviving it, points the same way: feature count correlates positively with
mean sentiment (\(\rho = +0.471\), \(p = .015\)) and mean score
(\(\rho = +0.422\), \(p = .032\)).

The lean and full groups also differ in overall rating, so bloat complaints
here might track general dissatisfaction rather than interface surfacing. Two
things speak against that: the complaint-rate gap survives dropping the two
smallest apps by volume, and the co-occurrence result below is estimated within
reviews rather than between apps, so it does not depend on the lean/full split
at all. We report the app-level comparison as exploratory and rest the
interface claim on the within-review evidence.

\begin{figure}[t]
  \centering
  \includegraphics[width=\linewidth]{fig08_rq5_bloat_curve}
  \caption{Bloat-complaint rate against feature count. The leanest apps draw bloat complaints at more than three times the rate of the fullest.}
  \label{fig:rq5-bloat-curve}
\end{figure}

\subsection{Bloat Complaints Are Interface Complaints}

If feature count is not the driver, what is? A review complaining about bloat
is 23.98 times more likely than base rate to also complain about the interface
(Fisher exact OR \(= 29.41\), \(p = 7.4\times10^{-137}\), \(n = 324{,}687\)).
The informative cell is small, and we give it rather than let the corpus-wide
\(n\) stand in: 134 reviews carry both complaints, 94 fusing them into a single
phrase and 40 keeping them in separate sentences. Restricting to those 40, the
lift is still 7.62\(\times\) (OR \(= 8.78\), \(p = 9.3\times10^{-24}\)). The
two lexicons share no keyword and no regex pattern, so a fused sentence carries
two distinct terms rather than one phrase counted twice. We report
7.62\(\times\) as the conservative bound but foreground the fused cases, since
they are literally what users say: bloat is voiced \emph{as} a property of the
interface.

\begin{quote}
\itshape ``Widget takes up too much space, [5\texttimes{}2], but displays
information only in [4\texttimes{}1]. Remaining blocks are blank, unusable for
other apps, just hogging space. Otherwise, good app, used daily. Has more than
timetables which for me is just unimportant clutter.''
\end{quote}
\hspace*{1em}--- Athan: Prayer Times \& Al Quran, 4 stars

\subsubsection{Provenance}
\label{sec:rq5-provenance}

The co-occurrence hypothesis and its test were specified on 2026-08-26, before
the full corpus was run, and that pre-registration stands. The decision to
foreground the fused cases, rather than lead with the conservative
separate-sentence bound, was made afterward with the co-occurrence figures and
labelled sub-topics already visible. The hypothesis and test were
pre-specified; the choice of reporting emphasis was not.

\subsection{The Sub-Topic Model Agrees, from a Different Direction}

A topic model fitted on \texttt{complexity\_bloat} reviews alone corroborates
the interface framing independently. Its labelled topics are dominated by
simplicity as something users \emph{praise} --- \emph{app praised for
simplicity} recurs across three topics, alongside \emph{praise for minimalistic
UI} and \emph{app praised for having no ads} --- with clutter, advertising, and
slowness (\emph{ads cause confusion while using app}, \emph{complaint for app
being slow}) on the negative side. Users are evaluating how much the interface
puts in front of them, not how many features the app contains. Topic structure
here is seed-sensitive (6--11 topics across five fixed seeds, mean pairwise ARI
0.452), so we report recurring themes rather than a fixed partition.

\subsection{Companion Hardware: A Case Study, \(N = 1\)}

One application, iQIBLA Life, ships a companion hardware device; of the 25
without one, 23 mention companion hardware only to request it. Inside iQIBLA
Life, hardware-aspect sentiment is 37.7\% negative and 45.9\% positive (518
mentions); across the 23 apps without hardware, mentions run 49.2\% negative
and 20.8\% positive (1,994 mentions), largely unmet requests for watch or ring
support rather than complaints about a device that exists. Shipping companion
hardware is therefore associated with \emph{better} hardware-topic sentiment
--- the reverse of a bloat story, and consistent with \S\ref{sec:rq3}'s finding
that unmet requests read more positively than usage reports of a feature that
can fail. The \texttt{companion\_hardware} tagger runs at 40\% precision (95\%
CI [12, 77]), so of 2,512 tagged mentions roughly 1,005 are genuine (interval
[295, 1{,}932]). With one treated application this is a descriptive case study;
no between-app estimate is possible.

\begin{quote}
\itshape ``app is not user friendly. often face problems of ring connectivity
and staff handling the problems is very incompetent and highly
unprofessional''
\end{quote}
\hspace*{1em}--- iQIBLA Life, 1 star

